% 电路与电子学III 课程笔记 —— 第6章（教材：孙立山《电路理论基础》第四版 第6章）
\chapter{非正弦周期电流电路}

\noindent
电气、电子及通信工程中广泛存在非正弦周期电流电路，其中的电流、电压是时间的非正弦周期函数。本章用\textbf{傅里叶级数}把非正弦周期激励分解为恒定分量与一系列不同频率的正弦分量，再用\textbf{叠加定理}把问题化为“直流电路＋若干正弦电路”的叠加问题；并讨论非正弦周期量的有效值、平均功率的计算，以及对称三相电路中的谐波。

\section{非正弦周期电流和电压}
非正弦电流电路是普遍存在的，其来源可归纳为四类：
\begin{enumerate}
  \item 电路中\textbf{同时有直流电源和正弦电源}作用时，电流一般既不是直流也不是正弦电流，而是非正弦周期电流。若电路是线性的，可由叠加定理分别计算各电源单独作用的响应，再把瞬时表达式相加；
  \item \textbf{电源本身}就是非正弦周期函数，例如方波、锯齿波电压源作用下的响应；
  \item 电路中含\textbf{非线性元件}，即使激励是正弦波，响应一般也是非正弦周期函数。典型例子是半波整流电路：输入正弦电压经二极管（单向导电）后，输出成为单一极性的非正弦周期电压；
  \item 实际电源并不理想：整流得到的直流电压仍存在周期性的\textbf{纹波}；发电机提供的电压也难以做到理想正弦波，存在一定的\textbf{畸变}。若必须考虑纹波或畸变的影响，就应建立非正弦周期电流电路模型。
\end{enumerate}

上述第 2、3 类情形都可以按第 1 类的思路处理：把非正弦周期激励分解为傅里叶级数，其中包含恒定分量和一系列不同频率的正弦分量，相当于有直流电源和多个不同频率的正弦电源同时作用于电路，再用线性电路的叠加定理计算。

\corepoint{非正弦周期电路的分析本质是“分解—分别计算—瞬时值叠加”。整章的三条主线就是：怎么分解（傅里叶级数）、怎么算有效值与功率、怎么按频率分别计算。}

\section{非正弦周期函数的傅里叶级数}
工程上遇到的周期函数 $f(t)$（周期 $T$，$\omega_1=2\pi/T$）总可以分解为傅里叶级数
\keyformula{f(t)=a_0+\sum_{k=1}^{\infty}\left[a_k\cos\left(k\omega_1t\right)+b_k\sin\left(k\omega_1t\right)\right]}
式中 $a_0$、$a_k$、$b_k$ 称为\textbf{傅里叶系数}。傅里叶级数收敛的充分条件是 $f(t)$ 满足狄里赫利条件：在一个周期内（1）连续或存在有限个第一类间断点；（2）只有有限个极值点。工程上遇到的非正弦周期函数一般都能满足。

各系数的计算（利用三角函数系的正交性）：
\keyformula{a_0=\frac{1}{T}\int_0^{T}f(t)\,\mathrm{d}t
=\frac{1}{2\pi}\int_0^{2\pi}f(t)\,\mathrm{d}\left(\omega_1t\right)}
\keyformula{a_k=\frac{2}{T}\int_0^{T}f(t)\cos\left(k\omega_1t\right)\mathrm{d}t,\qquad
b_k=\frac{2}{T}\int_0^{T}f(t)\sin\left(k\omega_1t\right)\mathrm{d}t}
其中 $a_0$ 就是 $f(t)$ 在一个周期内的\textbf{平均值}。

\subsection{振幅—初相形式}
把同频率的余弦项与正弦项合并，得到电工中更常用的形式
\keyformula{f(t)=A_0+\sum_{k=1}^{\infty}A_{mk}\cos\left(k\omega_1t+\psi_k\right)}
其中
\begin{equation}
A_{mk}=\sqrt{a_k^2+b_k^2},\qquad \psi_k=\arctan\frac{-b_k}{a_k}
\end{equation}
（$\psi_k$ 的象限要由 $a_k$、$b_k$ 的符号共同确定。）

\subsection{各项的物理含义}
从波形观点看：
\begin{itemize}
  \item $A_0=a_0$ 是常量，称为\textbf{恒定分量}（直流分量）；
  \item $A_{m1}\cos\left(\omega_1t+\psi_1\right)$ 的频率与 $f(t)$ 相同，称为\textbf{基波}（一次谐波），$A_{m1}$、$\psi_1$ 为基波的振幅与初相；
  \item $A_{mk}\cos\left(k\omega_1t+\psi_k\right)$ 的频率为基波的 $k$ 倍，称为 \textbf{$k$ 次谐波}。
\end{itemize}
由于傅里叶级数收敛，一般谐波次数越高、振幅越小。以谐波次数（或角频率）为横坐标、以振幅 $A_{mk}$ 为纵坐标画出的图形称为\textbf{振幅频谱}，图中竖线（谱线）的长度表示该次谐波的振幅，相邻两谱线的间隔等于基波角频率 $\omega_1$。这种谱线间具有固定间隔的频谱称为\textbf{离散频谱}。同理可画\textbf{相位频谱}，表示各次谐波初相 $\psi_k$ 随 $k\omega_1$ 变化的情形。把周期函数分解为恒定分量、基波和各次谐波的做法称为\textbf{谐波分析}。

\subsection{利用波形对称性简化计算}
实际波形的对称性可使许多系数为零，从而免去积分：
\begin{center}
\begin{tabular}{@{}lll@{}}
\toprule
波形性质 & 数学条件 & 结论 \\
\midrule
偶函数 & $f(-t)=f(t)$ & $b_k=0$，只含余弦项（及恒定分量）\\
奇函数 & $f(-t)=-f(t)$ & $a_0=0$、$a_k=0$，只含正弦项 \\
半波对称 & $f\left(t+T/2\right)=-f(t)$ & 只含奇次谐波，无恒定分量与偶次谐波 \\
\bottomrule
\end{tabular}
\end{center}

几个应记住的结论（$A$ 为方波幅值）：
\begin{itemize}
  \item 幅值为 $A$ 的对称方波（奇函数、半波对称）：$f(t)=\dfrac{4A}{\pi}\left[\sin\omega_1t+\dfrac{1}{3}\sin3\omega_1t+\dfrac{1}{5}\sin5\omega_1t+\cdots\right]$，即只含奇次谐波，振幅按 $1/k$ 下降；
  \item 三角波：只含奇次谐波，振幅按 $1/k^2$ 下降；
  \item 全波整流波形：只含恒定分量与偶次谐波。
\end{itemize}
若只测得波形曲线而无解析式，可把周期 $T$ 分为 $N$ 个相等时段，在 $N$ 个分点上取纵坐标，用求和公式近似代替上述积分，求出 $a_0$、$a_k$、$b_k$。

\corepoint{分解时先看对称性：奇函数只留正弦项、偶函数只留余弦项、半波对称只留奇次谐波。这一步能把积分量减半甚至减到四分之一，考试中省下的时间相当可观。}

\section{非正弦周期量的有效值、平均功率}
\subsection{有效值}
周期量的有效值等于其瞬时值的方均根值
\begin{equation}
I=\sqrt{\frac{1}{T}\int_0^{T}i^2(t)\,\mathrm{d}t}
\end{equation}
把 $i(t)$ 的傅里叶级数代入并展开平方，将得到四类积分：恒定分量平方的积分、同频率余弦平方的积分、不同频率乘积的积分（由三角函数的正交性为零）、恒定分量与余弦乘积的积分（为零）。结果
\keyformula{I=\sqrt{I_0^2+I_1^2+I_2^2+\cdots+I_k^2+\cdots}
=\sqrt{I_0^2+\sum_{k=1}^{\infty}I_k^2}}
式中 $I_0=a_0$ 为恒定分量，$I_k=A_{mk}/\sqrt{2}$ 为 $k$ 次谐波的有效值。即：

\textbf{任意周期量的有效值等于它的恒定分量、基波分量与各谐波分量有效值的平方和的平方根。}电压同理。

\subsection{平均功率}
设一端口网络的端口电压、电流取关联参考方向且为同频率的非正弦周期量，则平均功率是瞬时功率在一个周期内的平均值
\begin{equation}
P=\frac{1}{T}\int_0^{T}ui\,\mathrm{d}t
\end{equation}
把 $u$、$i$ 的傅里叶级数代入，由三角函数的正交性，只有\textbf{同频率}的电压与电流相乘的项积分不为零：
\begin{equation}
\frac{1}{T}\int_0^{T}U_{mk}\cos\left(k\omega_1t+\psi_{uk}\right)\cdot I_{mk}\cos\left(k\omega_1t+\psi_{ik}\right)\mathrm{d}t
=U_kI_k\cos\varphi_k
\end{equation}
其中 $U_k=U_{mk}/\sqrt{2}$、$I_k=I_{mk}/\sqrt{2}$，$\varphi_k=\psi_{uk}-\psi_{ik}$ 是 $k$ 次谐波电压超前于 $k$ 次谐波电流的相位差。于是
\keyformula{P=U_0I_0+\sum_{k=1}^{\infty}U_kI_k\cos\varphi_k
=P_0+P_1+P_2+\cdots}

结论：
\begin{itemize}
  \item 非正弦周期电流电路的平均功率等于\textbf{恒定分量、基波分量和各次谐波分量分别产生的平均功率之和}；
  \item \textbf{只有同频率的电压与电流才能产生平均功率}，不同频率的电压、电流相乘在一周期内的平均值为零；
  \item 因此可以“按频率分别算功率，再相加”，而不必做时域积分。
\end{itemize}

视在功率仍定义为 $S=UI$（$U$、$I$ 为非正弦量的有效值），此时一般 $S\neq\sqrt{P^2+Q^2}$（因含不同频率成分，存在畸变功率），故非正弦电路的功率因数 $\lambda=P/S$ 只是比值，不再等于 $\cos\varphi$。

\corepoint{有效值是“平方和开方”（各次谐波有效值平方相加），平均功率是“同频率配对相乘再相加”。千万别把不同次谐波的电压、电流相乘算成功率，也别把有效值直接相加。}

\section{非正弦周期电流电路的计算}
已知非正弦周期性激励和电路参数，可按三个步骤计算非正弦周期响应：

\begin{enumerate}
  \item \textbf{分解}：把给定的非正弦周期性激励分解为恒定分量和各正弦交流分量（查表或按傅里叶系数公式计算）；
  \item \textbf{分别计算}：分别计算电路在上述恒定分量和各交流分量单独作用下的响应。求恒定分量响应用直流电路的方法（电感短路、电容开路）；求各次交流分量的响应用正弦电流电路的相量法；
  \item \textbf{叠加}：把恒定分量和各交流分量的响应\textbf{瞬时表达式}相加，得到响应的时间函数。
\end{enumerate}

第 2 步中必须\textbf{按频率分别计算}，因为电感、电容对不同频率的谐波呈现不同的电抗：
\keyformula{X_{Lk}=k\omega_1L,\qquad X_{Ck}=-\frac{1}{k\omega_1C}}
即谐波次数越高，感抗越大、容抗越小。这正是滤波器（第 7 章）能分离不同频率成分的根据。

第 3 步的关键：表示各次交流的相量是代表\textbf{不同频率}正弦量的，它们对应不同相量方程的解，\textbf{不能把这些相量直接相加}（不同频率的相量相加没有意义）。必须先把各相量变换为瞬时表达式，再作瞬时值叠加。

求得响应的时间函数后，用上节公式可进一步求响应的有效值和平均功率。

\corepoint{三条铁律：（1）直流分量下电感短路、电容开路；（2）$k$ 次谐波下 $X_L$ 乘 $k$、$X_C$ 除 $k$；（3）叠加只能在\textbf{瞬时值}层面做，相量不能跨频率相加。}

\section{对称三相电路中的谐波}
对称三相非正弦周期电流通常为\textbf{上下半波镜像对称}，即 $f\left(t+T/2\right)=-f(t)$，此时电路中只含奇次谐波，且可按相序分为三组：

\begin{center}
\begin{tabular}{@{}lll@{}}
\toprule
谐波组 & 谐波次数 & 特点 \\
\midrule
正序组 & 1、7、13、$\cdots$（$6n+1$） & 三相互差 $120^\circ$，相序 A—B—C \\
负序组 & 5、11、17、$\cdots$（$6n-1$） & 三相互差 $120^\circ$，相序 A—C—B \\
零序组 & 3、9、15、$\cdots$（$3n$，$n$ 为奇数） & 三相\textbf{同相位}（相位差为零）\\
\bottomrule
\end{tabular}
\end{center}

分析方法：正序组和负序组谐波电路的分析方法与对称三相正弦电路完全一致（可取一相计算）；零序组谐波电路则要注意其\textbf{三相电压（电流）同相位}的特点，因而不能简单地“取一相”。

零序谐波（主要是三次谐波）的特殊影响：

\textbf{星形联结负载、无中性线时}：线电压、相电压和线电流中都\textbf{不存在零序谐波分量}。此时电源相电压中的零序分量与电源、负载两中性点间的电压相平衡，即在非正弦情况下，对称三相电路中电源与负载中性点之间的电压\textbf{不再等于零}，而等于电源电压中零序谐波电压的有效值
\begin{equation}
U_{N'N}=\sqrt{U_3^2+U_9^2+U_{15}^2+\cdots}
\end{equation}

\textbf{星形联结负载、有中性线时}：负载相电压和相电流中都含零序谐波，且三相零序电流同相位、都以中性线为返回路径，故中性线电流有效值为\textbf{一相零序谐波电流的 3 倍}
\begin{equation}
I_N=3\sqrt{I_{p3}^2+I_{p9}^2+\cdots}
\end{equation}
用单相计算电路求零序谐波响应时，为保证中性线阻抗上的电压符合实际，必须把\textbf{中性线阻抗乘以 3} 后画入单相电路（因为单相电路中只通过一相的三次谐波电流，而实际中性线中通过三相之和）。

\textbf{电源接成三角形时}：若电源相电压中含零序谐波（如三次谐波），三角形回路中三个相电压之和不再为零，而等于相电压中三次谐波电压有效值的\textbf{三倍}。这一电压可在三角形开口处用电压表测得。若把开口闭合，则不论外电路是否接通，三角形电源回路中始终存在三次谐波\textbf{循环电流}。要计算此电流必须修改电路模型，计入每相电源对三次谐波的内阻抗 $Z_3$，则
\begin{equation}
I_3=\frac{3U_3}{3Z_3}=\frac{U_3}{Z_3}
\end{equation}
此环流在每相内阻抗上的电压恰好等于电源相电压的三次谐波分量，二者互相抵消，故\textbf{线电压中不存在三次谐波分量}。其他零序谐波同理，因此
\begin{equation}
U_L=\sqrt{U_{L1}^2+U_{L5}^2+U_{L7}^2+\cdots}
\end{equation}
即线电压中不含零序谐波。为避免零序谐波循环电流、同时得到较高的线电压，\textbf{三相发电机一般接成星形而不接成三角形}。

\corepoint{三次谐波是三相电路里的“麻烦制造者”：三相同相位 → 星形无中性线时它顶在中性点之间，有中性线时它全部灌进中性线（$I_N=3I_{p3}$），三角形接法时它变成环流。记住这三条，谐波题基本就通了。}

\section*{本章小结}
\addcontentsline{toc}{section}{本章小结}
\begin{center}
\begin{tabular}{@{}p{0.22\textwidth}p{0.66\textwidth}@{}}
\toprule
项目 & 结论 \\
\midrule
傅里叶级数 & $f(t)=A_0+\sum A_{mk}\cos\left(k\omega_1t+\psi_k\right)$，$A_{mk}=\sqrt{a_k^2+b_k^2}$ \\
有效值 & $I=\sqrt{I_0^2+I_1^2+I_2^2+\cdots}$（平方和开方）\\
平均功率 & $P=U_0I_0+\sum U_kI_k\cos\varphi_k$（同频率才产生功率）\\
计算步骤 & 分解 → 按频率分别计算（直流法／相量法）→ 瞬时值叠加 \\
电抗与频率 & $X_{Lk}=k\omega_1L$，$X_{Ck}=-1/\left(k\omega_1C\right)$ \\
三相谐波分组 & 正序 $6n+1$、负序 $6n-1$、零序 $3n$（$n$ 为奇数）\\
零序谐波 & 线电压中不存在；有中性线时 $I_N=3I_{p3}$；三角形接法产生环流 \\
\bottomrule
\end{tabular}
\end{center}

\endinput
